\properlane{faithful-offering}{Mercy makes worship fruitful in a faithful life}

The person who enters confessing God's justice is later commanded to pay what has been vowed. Between those moments the Mass asks for purification, proclaims love and lordship, intercedes for a people, and enacts mysteries that heal. This movement places religious speech under the claim of fidelity. What is acknowledged, requested, and offered must reach conduct. At the same time, the final prayer asks for a remedy: a faithful life is the fruit of mercy, never a substitute offered in order to make mercy unnecessary.

Augustine's exposition of Psalm 75 supplies a demanding account of that movement. Confession establishes peace with God when the sinner ceases defending what God condemns. The renewed person who remembers the mercy already received does not treat renewal as a forgotten event. Vows then require payment, and payment requires God's help. Jerome's reading of Daniel adds intercession and perseverance in prayer; Thomas explains why the holy action has power to nourish the life it calls forth. Together they show how worship can become continuing fidelity.\footnote{Augustine, \textit{Psalms}, Ps.~75, paragraphs 3 and 9--11; Jerome, \textit{Daniel}, II, on 9:17--21; Thomas, \textit{Summa}, III, q.~79, aa.~1--2,6.}

\subsection{Confession releases the defense of self}
To acknowledge God's judgment is not merely to accept a general proposition that wrongdoing exists. Augustine's account of confession brings it home to the speaker. Peace with God begins when the person agrees with God's judgment against his own sin. This is not hatred of the person God created. It is refusal to identify that person so completely with his wrong that every criticism of the wrong must be treated as an attack on his existence. Mercy can then be sought as rescue rather than as vindication.\footnote{Augustine, \textit{Psalms}, Ps.~75, paragraph 3.}

The Introit puts that truth at the threshold. The servant acknowledges justice and then asks for mercy. He neither conceals the fault nor imagines that his fault gives him a claim to hopelessness. The Collect takes the next step by asking for release from corruption and a pure following of God alone. Repentance changes allegiance. Religious observance cannot become a protected compartment in which one praises God while leaving a cherished injustice untouched.

Paul's prisoner voice tests the sincerity of that change. Worthy walking takes the unglamorous forms of humility, meekness, patience, and mutual endurance. A person may prefer a conspicuous religious offering to the quiet task of bearing another's inconvenience. The Epistle makes the latter part of the former's truth. A gift laid before God does not authorize contempt for the person beside the giver. The shared Lord and baptism name a relation that persists beyond the moment of offering.

The Gradual places the giver within a prior gift. The people is chosen, the heavens are established by the Word, and their strength is the Spirit's work. Praise therefore begins in reception. In Bellarmine's account of inseparable Trinitarian action, the worshipper's offering does not initiate divine generosity. Augustine's spiritual heavens likewise stand through the Word. The offering returns a received life to its source.\footnote{Bellarmine, \textit{Psalms}, Ps.~32, commentary on vv.~6 and 12; Augustine, \textit{Psalms}, Ps.~32, first exposition, paragraph 6.}

\subsection{The Lord addresses the giver}
The Alleluia asks that prayer reach God before the Gospel is heard. That order gives listening a posture: the hearer is a suppliant before becoming an examiner of the words spoken. In the Gospel the lawyer asks the right kind of question about the law, yet the answer reaches the person who asks it. Chrysostom's account of his improvement offers the possibility that questioning can become conversion. Augustine's alternative account permits an initial testing that was cautious rather than malicious. Either way, truth is offered for reception rather than merely for victory in discussion.\footnote{Chrysostom, \textit{Matthew}, hom.~71; Augustine, \textit{De consensu evangelistarum}, II.73.141--142.}

Jesus' own question deepens that reception. Knowing which command is greatest does not yet exhaust the hearer's encounter with him. David's son is also David's Lord. The person who speaks about God is now addressed by the one to whom allegiance is owed. Chrysostom's Christological reasoning preserves both the human descent and the divine dignity: the problem lies in making the descent the entire answer. An offering made to this Lord cannot be a payment that leaves the giver otherwise self-owned.

The two love commands give that belonging its moral form. Worship is directed to God with the whole self, and the neighbor is loved as oneself. A strong act of devotion that leaves the neighbor deliberately excluded contradicts the command it honors. Conversely, service of the neighbor does not make the identity of the Lord irrelevant. The Gospel holds both together: who he is and what he asks determine the fidelity of the giver.

\subsection{Daniel brings the people's wound before God}
The Offertory turns from listening to intercession. Daniel's prayer asks for divine attention to sanctuary and people. The name of God upon them becomes the ground of appeal, while the full chapter rejects reliance on their own righteousness. The worshipper's approach thus includes a common wound. The act of offering is not an escape into a private account of innocence.

Jerome explains the sanctuary petition through promised restoration. Daniel seeks the fulfillment of God's word in the actual desolation before him. Divine faithfulness makes prayer possible, but it does not make prayer superfluous. The promise becomes the reason to ask. Jerome's treatment of God's eyes and ears similarly brings the language to its effect: the request is that God attend to the people's need, not that a bodily deficiency in God be repaired.\footnote{Jerome, \textit{Daniel}, II, on Dan.~9:17--18.}

At verse 20 Jerome offers the alternatives of personal confession as one of the people and humble identification with their guilt. Both make Daniel's solidarity active. He speaks for the people before God; he does not merely deliver a verdict upon them from a safe distance. At verse 21 Jerome interprets the mention of evening sacrifice through prayer persevering from morning to evening. The appeal continues; the holy person's attention is not exhausted by one eloquent sentence.\footnote{Jerome, \textit{Daniel}, II, on Dan.~9:20--21.}

That endurance has a moral analogue in the Epistle. To carry another in prayer while refusing every cost of patient charity would divide intercession from the good it seeks. Asking mercy for someone with whom one must also speak truthfully, then seeking a fitting time and manner for that speech, is one concrete form of it. Thomas's account of charitable correction allows both endurance and action. Prayer does not have to become an alternative to a necessary conversation, and the conversation does not have to become an exercise in punishment.\footnote{Thomas, \textit{Super Ephesios}, ch.~4, lect.~1.}

\subsection{The holy action heals the one who offers}
The Secret refuses to leave the meaning of offering at the level of sincere intention. It asks that the sacred things enacted free the people from offenses past and future. Thomas locates sacramental efficacy in Christ present and in his Passion, and explains its mode through nourishment. Food does more than remind a hungry person that eating would be desirable; it sustains life. So the sacramental action gives grace and strengthens the charity whose practice the Gospel commands.\footnote{Thomas, \textit{Summa}, III, q.~79, a.~1, corpus.}

The analogy also explains why preservation and subsequent freedom can coexist. Nourishment strengthens a living person without turning him into someone incapable of later harming himself. In Thomas's explicit account of free will, the future-directed petition is no advance absolution.

\Needspace{3\baselineskip}
The faithful ask to be guarded and strengthened, and their later choices remain consequential. The holy action reaches the future by giving life to be lived, not by making future conduct irrelevant.\footnote{Thomas, \textit{Summa}, III, q.~79, a.~6, corpus and ad 1.}

Augustine's discussion of remembrance in Psalm 75 gives continuity to that life. He relates baptismal renewal and the daily sacrificial remembrance of Christ to the need not to forget who has healed us. Repetition in worship can therefore sustain a real history of grace. It need not mean that an earlier gift was unreal, nor that gratitude may replace conversion. Remembering a mercy means allowing it to govern the life that follows.\footnote{Augustine, \textit{Psalms}, Ps.~75, paragraphs 9--10.}

\subsection{Vows are paid in ordinary fidelity}
At Communion the command to vow and pay brings speech to its test. Augustine distinguishes obligations common to Christians from particular vows freely undertaken. Ordinary fidelity is not reserved to people with a special religious commitment: refusal of hatred, pride, and dishonesty belongs to all. Particular commitments add their own demands, and the mere excellence of an undertaking does not excuse conduct that contradicts charity. His comparison of a humble married person and a proud consecrated virgin makes the point sharply: the higher profession does not make pride holy.\footnote{Augustine, \textit{Psalms}, Ps.~75, paragraph 11.}

Augustine's distinction imposes no new private vow on every hearer. The present question may be how to honor an existing responsibility, speak honestly, carry out a promised service, or stop using another's fault as permission for one's own. A more dramatic promise can sometimes distract from the duty already at hand. Augustine's warning concerns both making and fulfilling commitments; his encouragement rests on the help of the Lord to whom the commitment is made.

Before the feared sovereign of the final chant, this fidelity is more than an agreeable self-image. Bellarmine reads God's power over rulers concretely; Augustine finds pride stripped of its false sovereignty and urges mastery over one's own bodily conduct. Gifts do not purchase exemption from judgment. Those around the Lord offer in humility because the truth and the Lord at their center are common, not possessions their influence can control.\footnote{Bellarmine, \textit{Psalms}, Ps.~75, commentary on vv.~11--12; Augustine, \textit{Psalms}, Ps.~75, paragraphs 12--13.}

\subsection{The four senses of faithful offering}
\begin{fourSenses}
\item[Literal.] Confession, supplication, commands, vows, and prayers for healing retain their distinct purposes. In the Mass's actual sequence they accompany an approach to God, hearing, intercession, a holy action, reception, and a request for its lasting fruit. Daniel's sanctuary and the psalm's divine judgment give the sequence more than a private emotional meaning.

\item[Allegorical.] David's Lord is the saving Christ in whom the people approaches the Father. His Passion gives the mysteries their efficacy; his body carries the afflicted in a common voice. The Christian offering depends upon this prior divine gift.

\item[Moral.] Worship bears fruit when confession becomes renunciation of self-justification, intercession becomes patient care, and promised fidelity becomes conduct. Particular vows retain their seriousness, while the common law of love remains binding on every state of life. Divine help enables the response rather than excusing its absence.

\item[Anagogical.] The Lord before whom princes lose their power grants a remedy whose end is eternal. Thomas directs present sacramental nourishment toward perfect peace with the saints. The Postcommunion's final request therefore reaches beyond a better week to the completion of a life in faithful communion with God.\footnote{Thomas, \textit{Summa}, III, q.~79, a.~2, corpus.}
\end{fourSenses}
